\subsection{Limitation}
 
While our method significantly enhances the performance of codecs for LLMs by integrating semantic information, it does come with certain trade-offs. According to the principle of "no free lunch," improving one aspect of a system often involves compromises in others. In the case of our enhanced codecs, the primary limitation lies in their potential impact on the original functionality of codecs, which is compression for information transmission. The introduction of a semantic extraction layer adds additional computational overhead, potentially increasing the time required for processing. This can affect the efficiency of the codec when used in applications where rapid data compression and transmission are critical. Consequently, while our approach offers substantial benefits for semantic understanding and audio processing, it may not be as effective in contexts where high-speed data compression is paramount.
\begin{table}[h!]
\centering
\resizebox{0.35\textwidth}{!}{%
\begin{tabular}{lccc}
\hline
\textbf{Model} & \textbf{Mel DT.} & \textbf{STFT DT.} & \textbf{UTMOS} \\
\hline
Baseline ($n_q$=$1$) & 0.79 & 0.73 & 2.96 \\
Baseline ($n_q$=$8$) & 0.54 & 0.58 & 3.72 \\
X-codec ($n_q$=$1$) & 0.86  & 0.77  & 3.71 \\
X-codec ($n_q$=$8$) & 0.62 & 0.63  & 4.01 \\
\hline
\end{tabular}
}
\caption{Comparison of reconstruction based on Mel DT., STFT DT., and UTMOS metrics using 1000 LibriTTS speech samples. “DT.” is short for distance}
\end{table}

Furthermore, the integration of semantic layers can slightly impair certain acoustic metrics such as Mel and STFT distance, which are crucial for maintaining the fidelity of compressed audio. However, it is essential to note that these trade-offs are counterbalanced by significant improvements in human auditory perception, as evidenced by the UTMOS scores.

 
 
 
 
\section{Conclusion}

In this paper, we introduced X-codec, an advanced audio codec that integrates semantic information through self-supervised learning models to enhance performance in large language models, specifically in text-to-speech synthesis, music continuation, and general audio classification tasks. Our evaluations demonstrate that X-codec significantly improves semantic understanding and audio generation quality across a variety of domains.  

% In this paper, we first introduce the a unified audio tokenization approach that elegantly integrates semantic and acoustic information and can be used in language foundation model in a plug-and-play way. Then, our empirical studies demonstrates enhanced capabilities in both semantic preservation and audio quality, which is beneficial in understanding and generation tasks. Moreover, the exploration of the application potentials of X-Codec in GPT-style audio generation, showing promising and comparative results. With our futher released code, checkpoints or even API, our novel idea and solid tokenization method have great potential to help the community to build the promising and powerful audio foundation model in the future.